\subsubsection{Symmetry Boundary Condition}
\label{SEC:BC:Numerical:Symmetry}

This boundary condition is generally used to reduce the size of the overall
computational domain by taking advantage of flow symmetry for a given problem.
%
For example, instead of running a simulation of a full airplane, a symmetry
boundary can be created along the centerline of the airplane geometry so the
simulation is only performed on half of the airplane.
%
Although intended for different uses, this boundary condition is equivalent to
the inviscid\slash slip wall boundary condition from \sref{SEC:BC:Wall:Slip}.

%%%%%%%%%%%%%%%%%%%%%%%%%%%%%%%%%%%%%%%%%%%%%%%%%%%%%%%%%%%%%%%%%%%%%%%%%%%%%%%%

\NOTE It should be warned that using symmetry boundary conditions to take
advantage of flow symmetry is generally not recommended for quality LES.
%
The mean flow may be symmetric over long time-averages, but the instantaneous
turbulent flow structures are \emph{not} symmetric.

%###############################################################################

